\section{Verification}\label{sec:verif}

\subsection{Method}
Every reference value in this section is analytic or comes from an independent computation (for example numerical quadrature of the elastica integral); none comes from \fea{} itself. The benchmarks are the cases
of \texttt{fea\_engine/tests/test\_benchmarks\_convergence.py}, and the scripts in \texttt{paper/experiments/} import that test module's helpers, so the numbers below are what the test suite checks.
Convergence orders are the slope of the error against element size $h$ on a log--log scale: \texttt{fitted} uses all refinements and \texttt{last two} the two finest. Results below were produced by \texttt{paper/experiments/run\_all.py} on the machine in \cref{sec:perf}.
Verification here means agreement with a reference on the stated cases; it does not establish correctness outside them.

\subsection{Finite element benchmarks}
\paragraph{Patch tests.}
On distorted meshes, a linear displacement field must be reproduced exactly by every continuum element and the recovered stress must be constant. \Cref{tab:patch} shows that all eight continuum elements pass to rounding error; the largest displacement or stress error is \patchMaxErr{}.

\begin{table}[t]
\centering\small
\caption{Patch test on distorted meshes: relative error in displacement and in recovered stress.}\label{tab:patch}
\begin{tabular}{lrr}
\toprule
Element & displacement error & stress error \\
\midrule
\texttt{Quad4PlaneStress} & $5.5\times10^{-16}$ & $5.2\times10^{-15}$ \\
\texttt{Quad8PlaneStress} & $1.0\times10^{-15}$ & $7.6\times10^{-15}$ \\
\texttt{Tri3PlaneStress} & $3.7\times10^{-16}$ & $4.9\times10^{-15}$ \\
\texttt{Tri6PlaneStress} & $1.8\times10^{-15}$ & $1.7\times10^{-14}$ \\
\texttt{Hex8Solid3D} & $1.8\times10^{-16}$ & $8.9\times10^{-16}$ \\
\texttt{Hex20Solid3D} & $1.2\times10^{-15}$ & $2.7\times10^{-15}$ \\
\texttt{Tet4Solid3D} & $3.6\times10^{-16}$ & $8.9\times10^{-16}$ \\
\texttt{Tet10Solid3D} & $1.2\times10^{-15}$ & $1.9\times10^{-15}$ \\
\bottomrule
\end{tabular}

\end{table}

\paragraph{Pure bending.}
A plane-stress beam under an end moment has an exact elasticity solution with quadratic displacement. The quadratic elements (Quad8, Tri6) contain that solution and reproduce the tip deflection to rounding error (largest deviation of the ratio from one: \quadBendingMaxErr{}).
The linear elements converge at the theoretical second order but carry the known bending-locking offset: on the finest mesh the tip deflection is \quadBendFine{} of exact for Quad4 and \triBendFine{} for Tri3 (\cref{tab:orders}). The Tri3 order measured on the last two refinements is closer to two than the fitted order over all three, consistent with a pre-asymptotic range.

\paragraph{Beam and bar vibration.}
The first frequency of an Euler--Bernoulli cantilever (exact value \ebFirstHz{}~Hz for the chosen section and length) converges at fourth order in the number of Hermite elements, as expected for the cubic interpolation; with 20 elements the relative error of the first frequency is \ebRelErrTwenty{}, and the error grows for higher modes. The first axial frequency of a fixed-free bar of Quad4 elements converges at second order.

\paragraph{Shear-deformable beam.}
The Reissner beam converges at second order to the Timoshenko closed-form tip deflection $PL^3/(3EI)+PL/(GA\kappa)$, with a relative error of \timoRelErr{} on the finest of \timoN{} refinements.

\paragraph{Large deflection.}
For a cantilever with an end load $PL^2/EI=1$, the tip deflection and the end shortening from the nonlinear Reissner beam (10 load steps) are compared with the elastica integrated by quadrature (\cref{tab:elastica}). With 40 elements the error in the vertical deflection is \elasticaWerr{} and in the shortening \elasticaShortErr{}. The deflection error does not fall monotonically at a fixed rate between the three meshes (the value changes sign relative to the reference between 20 and 40 elements), so we do not report an order for this case.

\begin{table}[t]
\centering\small
\caption{Cantilever elastica, $PL^2/EI=1$: tip deflection $w$ and end shortening, both divided by $L$, and relative errors against the quadrature reference.}\label{tab:elastica}
\begin{tabular}{rrrrr}
\toprule
elements & $w/L$ & shortening$/L$ & error in $w$ & error in shortening \\
\midrule
10 & 0.301166 & 0.056141 & $1.8\times10^{-3}$ & $5.2\times10^{-3}$ \\
20 & 0.301611 & 0.056365 & $3.6\times10^{-4}$ & $1.2\times10^{-3}$ \\
40 & 0.301722 & 0.056421 & $3.5\times10^{-6}$ & $2.2\times10^{-4}$ \\
reference & 0.301721 & 0.056433 & -- & -- \\
\bottomrule
\end{tabular}

\end{table}

\begin{table}[t]
\centering\small
\caption{Observed convergence orders. Plane-stress rows use the last two refinements; other rows use the fit over all refinements.}\label{tab:orders}
\begin{tabular}{lrr}
\toprule
Benchmark & observed order & theory \\
\midrule
Plane-stress pure bending, Quad4 & 1.96 & 2 \\
Plane-stress pure bending, Tri3 & 1.79 & 2 \\
Euler--Bernoulli beam, first frequency & 3.98 & 4 \\
Fixed-free bar, first frequency (Quad4) & 2.00 & 2 \\
Reissner beam, Timoshenko tip deflection & 2.00 & 2 \\
\bottomrule
\end{tabular}

\end{table}

\begin{figure}[t]
\centering
\includegraphics[width=\textwidth]{figures/e5_verification.pdf}
\caption{Error against element size $h$ for pure bending (plane-stress, linear elements), the first frequencies, and the Timoshenko tip deflection. Dotted lines show the theoretical slopes.}\label{fig:verif}
\end{figure}

\subsection{Reduced-order model checks}
\paragraph{Parametric POD--Galerkin model.}
The model is a cantilever of 120 beam elements ($\eOneFull{}$ DOF) with two material regions whose bending rigidities $EI_1,EI_2$ are independent parameters drawn from $[0.5,5]$ in the units of the shipped example (the response is linear in the load, so relative errors do not depend on the scale). The ROM uses an affine decomposition, a basis of \eOneReduced{} POD modes built from \eOneNtrain{} full-order solves, and is tested on \eOneNtest{} held-out random parameter pairs.
The maximum relative tip-deflection error is \eOneErr{} and the mean is \eOneMeanErr{} (\cref{fig:e1}). The test parameters differ from the training ones, so this is an out-of-sample error.

\paragraph{Craig--Bampton synthesis.}
A clamped--clamped Euler--Bernoulli beam of \cmsFullSize{} free DOF is split in two halves joined at one node. The reference is the generalised eigenproblem of the unreduced beam (first frequency \cmsFirstHz{}~Hz). \Cref{tab:cms} gives the largest relative error over the first four frequencies against the number of fixed-interface modes kept per half.
With two modes per half (reduced size \cmsSizeTwo{}) the error is \cmsErrTwo{}, and with five (size \cmsSizeFive{}) it is \cmsErrFive{}; the error decreases steadily with more modes. Keeping all modes (\cmsSizeAll{} reduced DOF, equal to the full size) gives \cmsErrAll{}, which only checks that coupling and expansion are exact when nothing is truncated. Our error metric is the maximum relative frequency error and differs from the normalisation used in the package README.

\begin{table}[t]
\centering\small
\caption{Craig--Bampton error of the first four frequencies against reduced size.}\label{tab:cms}
\begin{tabular}{rrr}
\toprule
modes per half & reduced size & max rel. freq. error \\
\midrule
1 & 4 & $4.8\times10^{-1}$ \\
2 & 6 & $3.2\times10^{-3}$ \\
3 & 8 & $7.0\times10^{-4}$ \\
5 & 12 & $9.3\times10^{-5}$ \\
8 & 18 & $1.2\times10^{-5}$ \\
10 & 22 & $4.3\times10^{-6}$ \\
15 & 32 & $5.3\times10^{-7}$ \\
20 & 42 & $1.0\times10^{-7}$ \\
all & 78 & $5.2\times10^{-11}$ \\
\bottomrule
\end{tabular}

\end{table}

\begin{figure}[t]
\centering
\includegraphics[width=0.48\textwidth]{figures/e1_rom_speedup.pdf}\hfill
\includegraphics[width=0.48\textwidth]{figures/e3_cms.pdf}
\caption{Left: tip deflection of the two-region cantilever against the stiffness ratio, full model (markers) and ROM (lines). Right: Craig--Bampton frequency error against modes per half.}\label{fig:e1}
\end{figure}

\paragraph{Hyper-reduction.}
The accuracy of ECSW on held-out states and along a time trajectory is reported with its cost in \cref{tab:ecsw}.

\subsection{What this section does not show}
The benchmarks cover the element types and analyses named above on simple geometries. They do not cover three-dimensional plasticity, contact, shells or the Koiter--Newton solvers against independent references, nor error bounds of the certified type; those components are covered by the unit tests in the repository but are not evaluated here.
